% GENERATED FILE. Edit canonical lessons, then run npm run generate:books.
% canonical-source-hash: fnv1a64:3ec39629e00acdae
% canonical-lessons: SA-C153-itah, SA-C153-tatah, SA-C153-ubhayatah, SA-C153-paritah, SA-C153-samantatah

\chapter{From here, From there, On both sides, All around, On all sides}
\label{ch:sa-from-here-from-there-on-both-sides-all-around-on-all-sides}
\hlchaptermodality{153}

\begin{chapteropening}
\textbf{By the end of this chapter:} \emph{I can say from here, from there, on both sides, all around and on all sides.}

Complete the last lesson of chapter 153: I can say from here, from there, on both sides, all around and on all sides.
\end{chapteropening}

\section[\texorpdfstring{\sk{इतः} (itaḥ)}{इतः (itaḥ)}]{\sk{इतः} (itaḥ) — from here}

\label{lesson:SA-C153-itah}



\begin{quote}
Before the new one: say the Sanskrit for a calendar date, then the Sanskrit for a century.
\end{quote}

\subsection*{You'll want to know: \sk{इतः}}
\textbf{\sk{इतः}} — \emph{itaḥ} — \textquotedblleft{}from here\textquotedblright{}.

\begin{grammarlens}[title={What you've built}]
One more word for this chapter.
\end{grammarlens}

\subsection*{Your turn}
\begin{itemize}
\raggedright
  \item \emph{Say it:} \emph{itaḥ}
  \item \emph{Say it:} \emph{itaḥ}, once more
  \item \emph{Say it:} say \emph{itaḥ}
  \item \emph{Recall:} say the Sanskrit for a calendar date, then the Sanskrit for a century, then say \emph{itaḥ} again
\end{itemize}

\begin{tcolorbox}[breakable,colback=teal!4,colframe=teal!35!black,title={Before you move on}]
What does \sk{इतः} mean? (\textquotedblleft{}From here\textquotedblright{}.) Say it once more.
\end{tcolorbox}
\section[\texorpdfstring{\sk{ततः} (tataḥ)}{ततः (tataḥ)}]{\sk{ततः} (tataḥ) — from there, then}

\label{lesson:SA-C153-tatah}



\begin{quote}
Before the new one: say the Sanskrit for a century, then the Sanskrit for from here.
\end{quote}

\subsection*{You'll want to know: \sk{ततः}}
\textbf{\sk{ततः}} — \emph{tataḥ} — \textquotedblleft{}from there, then\textquotedblright{}.

\begin{grammarlens}[title={What you've built}]
One more word for this chapter.
\end{grammarlens}

\subsection*{Your turn}
\begin{itemize}
\raggedright
  \item \emph{Say it:} \emph{tataḥ}
  \item \emph{Say it:} \emph{tataḥ}, once more
  \item \emph{Say it:} say \emph{tataḥ}
  \item \emph{Recall:} say the Sanskrit for a century, then the Sanskrit for from here, then say \emph{tataḥ} again
\end{itemize}

\begin{tcolorbox}[breakable,colback=teal!4,colframe=teal!35!black,title={Before you move on}]
What does \sk{ततः} mean? (\textquotedblleft{}From there, then\textquotedblright{}.) Say it once more.
\end{tcolorbox}
\section[\texorpdfstring{\sk{उभयतः} (ubhayataḥ)}{उभयतः (ubhayataḥ)}]{\sk{उभयतः} (ubhayataḥ) — on both sides}

\label{lesson:SA-C153-ubhayatah}



\begin{quote}
Before the new one: say the Sanskrit for from here, then the Sanskrit for from there.
\end{quote}

\subsection*{You'll want to know: \sk{उभयतः}}
\textbf{\sk{उभयतः}} — \emph{ubhayataḥ} — \textquotedblleft{}on both sides\textquotedblright{}.

\begin{grammarlens}[title={What you've built}]
One more word for this chapter.
\end{grammarlens}

\subsection*{Your turn}
\begin{itemize}
\raggedright
  \item \emph{Say it:} \emph{ubhayataḥ}
  \item \emph{Say it:} \emph{ubhayataḥ}, once more
  \item \emph{Say it:} say \emph{ubhayataḥ}
  \item \emph{Recall:} say the Sanskrit for from here, then the Sanskrit for from there, then say \emph{ubhayataḥ} again
\end{itemize}

\begin{tcolorbox}[breakable,colback=teal!4,colframe=teal!35!black,title={Before you move on}]
What does \sk{उभयतः} mean? (\textquotedblleft{}On both sides\textquotedblright{}.) Say it once more.
\end{tcolorbox}
\section[\texorpdfstring{\sk{परितः} (paritaḥ)}{परितः (paritaḥ)}]{\sk{परितः} (paritaḥ) — all around}

\label{lesson:SA-C153-paritah}



\begin{quote}
Before the new one: say the Sanskrit for from there, then the Sanskrit for on both sides.
\end{quote}

\subsection*{You'll want to know: \sk{परितः}}
\textbf{\sk{परितः}} — \emph{paritaḥ} — \textquotedblleft{}all around\textquotedblright{}.

\begin{grammarlens}[title={What you've built}]
One more word for this chapter.
\end{grammarlens}

\subsection*{Your turn}
\begin{itemize}
\raggedright
  \item \emph{Say it:} \emph{paritaḥ}
  \item \emph{Say it:} \emph{paritaḥ}, once more
  \item \emph{Say it:} say \emph{paritaḥ}
  \item \emph{Recall:} say the Sanskrit for from there, then the Sanskrit for on both sides, then say \emph{paritaḥ} again
\end{itemize}

\begin{tcolorbox}[breakable,colback=teal!4,colframe=teal!35!black,title={Before you move on}]
What does \sk{परितः} mean? (\textquotedblleft{}All around\textquotedblright{}.) Say it once more.
\end{tcolorbox}
\section[\texorpdfstring{\sk{समन्ततः} (samantataḥ)}{समन्ततः (samantataḥ)}]{\sk{समन्ततः} (samantataḥ) — on all sides}

\label{lesson:SA-C153-samantatah}



\begin{quote}
Before the new one: say the Sanskrit for on both sides, then the Sanskrit for all around.
\end{quote}

\subsection*{You'll want to know: \sk{समन्ततः}}
\textbf{\sk{समन्ततः}} — \emph{samantataḥ} — \textquotedblleft{}on all sides\textquotedblright{}.

\begin{grammarlens}[title={What you've built}]
One more word for this chapter.
\end{grammarlens}

\subsection*{Your turn}
\begin{itemize}
\raggedright
  \item \emph{Say it:} \emph{samantataḥ}
  \item \emph{Say it:} \emph{samantataḥ}, once more
  \item \emph{Say it:} say \emph{samantataḥ}
  \item \emph{Recall:} say the Sanskrit for on both sides, then the Sanskrit for all around, then say \emph{samantataḥ} again
\end{itemize}

\begin{tcolorbox}[breakable,colback=teal!4,colframe=teal!35!black,title={Before you move on}]
What does \sk{समन्ततः} mean? (\textquotedblleft{}On all sides\textquotedblright{}.) Say it once more.
\end{tcolorbox}
